\section{Comparing products of shifted primes}\label{sec:layers}

We next bound the number of ways that two ordered lists of primes can
give the same product of their shifts.  The largest prime factors of the
shifts lie in separated intervals.  We expose the factors one interval
at a time, starting with the smallest interval.  At each stage only two
or three simultaneous prime conditions are needed.  This keeps the
constants uniform when the number of intervals grows.

The layer decomposition and the comparison of dual factorizations
go back to Maier and Pomerance \cite[Section~3]{MaierPomerance1988}
and were developed further in Ford's proof of Lemma~5.1
\cite{FordTotients2013}. We prove the version needed here,
including its dependence on the number of layers.  In particular, we
assume squarefreeness of the whole product above the last cutoff;
squarefreeness of each individual shift would not suffice for our
factor-allocation estimate.

\subsection{A uniform sieve for two or three forms}

\begin{lemma}\label{lem:sieve}
There are absolute constants $C$ and $y_0$ with the following property.
Let $y\ge y_0$, put $B_y=\log_2 y$, and let $1\le A,A'\le y$ be
integers.  Consider either the two forms
\[
 n,\quad An+1,
\]
or the three forms
\[
 n,\quad An+1,\quad A'n+1,\qquad A\ne A'.
\]
Write $k=2$ or $3$, respectively, and let $\mathcal P$ be the set of
positive integers at which all the indicated forms are prime.  Then,
uniformly for $X\ge3$,
\begin{equation}\label{eq:sieve-count}
 \#(\mathcal P\cap[1,X])
 \le C\frac{X B_y^k}{(\log X)^k}.
\end{equation}
Moreover, uniformly for $3\le U\le Y$,
\begin{equation}\label{eq:sieve-reciprocal}
 \sum_{\substack{n\in\mathcal P\\U\le n\le Y}}\frac1n
 \le C\frac{B_y^k}{(\log U)^{k-1}}.
\end{equation}
The coefficients need not be bounded in terms of $X$ or $U$.
\end{lemma}

\begin{proof}
For primitive, distinct, admissible integral linear forms
$a_jn+b_j$ with positive leading coefficients and nonzero determinant
$\Delta$ as defined below, Ford's uniform upper sieve
\cite[Theorem~2.5]{FordSieve2023} gives
\begin{equation}\label{eq:sieve-determinant}
 \begin{gathered}
 \#\{1\le n\le X:a_jn+b_j\text{ prime for every }j\}
 \ll_k\frac{X}{(\log X)^k}
 \left(\frac{\Delta}{\varphi(\Delta)}\right)^k,\\
 \Delta=\left|\prod_j a_j
       \prod_{j<l}(a_jb_l-a_lb_j)\right|.
 \end{gathered}
\end{equation}
Here admissibility means that no prime divides the product of the
forms for every integer input; the implied constant depends only
on $k$, including when the coefficients vary.  For our forms,
primitivity is automatic, and the determinants are $A$ and
$AA'(A-A')$, up to sign.  They are nonzero and have absolute value at
most $y^3$.

If our forms are inadmissible, an offending prime is at most $k$, since
each primitive linear form has at most one root modulo that prime.
At every simultaneous-prime input one of the forms must equal this
prime.  Each such equality has at most one solution, so there are at
most $k$ inputs.  Their contribution is absorbed in
\eqref{eq:sieve-count}.

For admissible forms, the elementary estimate
\[
 \frac{d}{\varphi(d)}\ll\log_2(d+10)\qquad(d\ge1)
\]
and \eqref{eq:sieve-determinant} prove \eqref{eq:sieve-count}.
To recall its uniformity, split the product over prime divisors of
$d$ at $\log(d+10)$.  Mertens' product estimate bounds the first part
by $O(\log_2(d+10))$.  For the remaining part, use
$\sum_{p\mid d}\log p\le\log d$ to bound the sum of reciprocal
prime divisors, giving a bounded factor.  Since $k$ takes only the
values $2$ and $3$, all constants are absolute.

Finally, partial summation of \eqref{eq:sieve-count}, with the
coefficients fixed, gives
\[
 \sum_{\substack{n\in\mathcal P\\U\le n\le Y}}\frac1n
 \ll B_y^k\left(\frac1{(\log Y)^k}
       +\int_U^Y\frac{\dd t}{t(\log t)^k}\right)
 \ll\frac{B_y^k}{(\log U)^{k-1}}.
\]
This also includes a possible prime at the lower endpoint.
\end{proof}

\subsection{The comparison estimate}

For a positive integer $n$ and a real number $Z$, its part supported
above $Z$ means $\prod_{p^a\parallel n,\ p>Z}p^a$.  We use
$S$-normality as defined in Definition~\ref{def:normal}.

The key term in the bound below is the exponent of $\log y$.
In its application, a strict simplex inequality makes
the weighted cutoff sum plus its error smaller than $1$ by a
controlled amount.
The exponent is consequently smaller than $-1$; that extra saving
absorbs the other factors, even after recovering the tuple data hidden
in the smooth remainder. We need an estimate for ordered pairs with
multiplicity: a bound only on the number of values admitting a collision
would not control how many tuples lie over those values.

\begin{proposition}\label{prop:comparison}
There are absolute constants $C$ and $y_0$ such that the following
holds.  Let $y\ge y_0$, put $B_y=\log_2 y$, and let $b$ be an integer
with $1\le b\le B_y$.  Suppose that real cutoffs satisfy
\[
 Y_0=y,\qquad \e^\e\le S\le Y_b,
 \qquad Y_{j+1}<U_j<Y_j\quad(0\le j<b).
\]
Put
\[
 \nu_j=\frac{\log_2 Y_j}{B_y}\quad(0\le j\le b),
 \qquad
 \mu_j=\frac{\log_2 U_j}{B_y}\quad(0\le j<b),
 \qquad
 \delta=\sqrt{\frac{\log_2 S}{B_y}}.
\]
Fix a positive integer $D$ with $P^+(D)\le Y_b$ and a real $r\ge1$.
Let $N$ count the ordered tuples
\[
 (p_0,\ldots,p_{b-1},q_0,\ldots,q_{b-1},D')
\]
in which $D'$ is a positive integer and the following conditions hold:
\begin{enumerate}[label=\textup{(\roman*)}]
 \item The $p_j,q_j$ are $S$-normal primes, $p_j\ne q_j$, and
 \[
  U_j\le P^+(p_j-1),P^+(q_j-1)\le Y_j
  \qquad(0\le j<b).
 \]
 \item $P^+(D')\le Y_b$ and
 \[
  D\prod_{j=0}^{b-1}(p_j-1)
  =D'\prod_{j=0}^{b-1}(q_j-1)\le\frac yr.
 \]
 \item The common product in \textup{(ii)} is squarefree on primes
 greater than $Y_b$.
 \item The part of $p_0-1$ supported above $Y_1$ exceeds $\sqrt y$.
\end{enumerate}
Then
\begin{equation}\label{eq:comparison}
 \begin{split}
 N\le{}&\frac{y}{rD}(CB_y^6)^b(b+1)^{\Omega(D)}
       (\log Y_b)^{1+3b\log(b(b+1))}\\
 &\quad\times
 (\log y)^{-2+\sum_{j=1}^{b-1}a_j\nu_j+E_b},
 \end{split}
\end{equation}
where
\begin{equation}\label{eq:comparison-error}
 E_b=2\sum_{j=1}^{b-1}(\nu_j-\mu_j)
      +\delta\sum_{i=2}^b(i\log i+i).
\end{equation}
Here $a_j$ is as defined in \eqref{eq:aj}.  Empty sums are zero;
in particular, when $b=1$ the cutoff in
\textup{(iv)} is $Y_1=Y_b$.
\end{proposition}

\begin{proof}
We first describe the layers and the weighted summation that combines
them.  We then count the top layer, estimate a general lower layer,
and finally sum over the small prime factors.

\paragraph{The layer configurations.}
For $0\le i\le b$ and $0\le j<b$, let
\[
 s_{i,j}=\prod_{\substack{p^a\parallel p_j-1\\p\le Y_i}}p^a,
 \qquad
 s'_{i,j}=\prod_{\substack{p^a\parallel q_j-1\\p\le Y_i}}p^a.
\]
Restricting the equality in condition~\textup{(ii)} to primes at most
$Y_i$ gives the common product
\[
 s_i^\#=D\prod_{j=0}^{b-1}s_{i,j}
        =D'\prod_{j=0}^{b-1}s'_{i,j}.
\]
For $1\le i\le b$, put
\[
 t_{i,j}=\frac{s_{i-1,j}}{s_{i,j}},\qquad
 t'_{i,j}=\frac{s'_{i-1,j}}{s'_{i,j}},\qquad
 t_i^\#=\frac{s_{i-1}^\#}{s_i^\#}.
\]
Each $t_i^\#$ is squarefree and supported on $(Y_i,Y_{i-1}]$.
For $j\ge i$ the factors $t_{i,j}$ and $t'_{i,j}$ equal $1$, so
\begin{equation}\label{eq:layer-two-allocations}
 t_i^\#=\prod_{j=0}^{i-1}t_{i,j}
        =\prod_{j=0}^{i-1}t'_{i,j}.
\end{equation}
Thus a layer consists of two allocations of the same squarefree
integer among $i$ labeled slots.

Let $\mathcal C_i$ be the set of configurations
\[
 \sigma_i=(D',s_{i,0},\ldots,s_{i,b-1},
                   s'_{i,0},\ldots,s'_{i,b-1})
\]
arising from tuples counted by $N$.  Restriction of prime factors
maps $\mathcal C_{i-1}$ into $\mathcal C_i$; the two allocations in
\eqref{eq:layer-two-allocations} determine each extension uniquely.
Reconstructing a tuple from its small prime factors runs from
$\mathcal C_b$ back to $\mathcal C_0$; the weighted count below
sums in the opposite direction, starting with the top layer.
The original tuples are in bijection with $\mathcal C_0$, since
$p_j=s_{0,j}+1$ and $q_j=s'_{0,j}+1$.  For $i\ge2$ we will bound,
uniformly in $\sigma_i$, the reciprocal extension sum
\[
 \sum_{\substack{\sigma_{i-1}\in\mathcal C_{i-1}\\
                  \sigma_{i-1}\text{ restricts to }\sigma_i}}
       \frac1{t_i^\#}.
\]
This is the appropriate weight because
$1/s_{i-1}^\#=1/(s_i^\#t_i^\#)$.

Every configuration considered here extends to a genuine solution.
In particular, its partial common product is at most $y/r$, each
shift is at most $y$, and every fixed or conditional coefficient
used below is at most $y$.  When enlarging a sum for an upper bound,
we retain any necessary restriction on a conditional coefficient
until the corresponding sieve estimate has been applied.

\paragraph{The top layer.}
Only the zeroth shift can have a prime factor greater than $Y_1$.
Consequently $t_1^\#$ is the same part of $p_0-1$ and $q_0-1$,
and $t_1^\#>\sqrt y$.  Normality, applied with endpoint $p_0-1$,
gives
\[
 \Omega(p_0-1)
 \le B_y+\log_2 S+\sqrt{B_y\log_2 S}\le3B_y.
\]
Write $t_1^\#=t'Q$, with $Q=P^+(t_1^\#)$.  The weaker bound
$\Omega(t_1^\#)\le4B_y$ already gives
\[
 Q\ge y^{1/(8B_y)}.
\]
After $\sigma_1$ and $t'$ are fixed, the completed primes have the
forms $AQ+1,A'Q+1$.  Their inequality implies $A\ne A'$, which
is exactly the nonzero determinant condition for the three-form
sieve.  With
\[
 X=\frac{y}{r s_1^\#t'},
\]
a nonempty extension has $X\ge Q\ge y^{1/(8B_y)}$.
Lemma~\ref{lem:sieve} therefore bounds the number of $Q$ by
\[
 \ll\frac{X B_y^3}{(\log X)^3}
 \ll\frac{y}{r s_1^\#t'}\frac{B_y^6}{(\log y)^3}.
\]
The integer $t'$ is squarefree and supported on $(Y_1,y]$.
Mertens' estimate gives
\[
 \sum_{t'}\frac1{t'}
 \le\prod_{Y_1<p\le y}\left(1+\frac1p\right)
 \ll\frac{\log y}{\log Y_1}.
\]
It follows that, for each $\sigma_1$,
\begin{equation}\label{eq:top-layer}
 \#\{\sigma_0\text{ restricting to }\sigma_1\}
 \le CB_y^6\frac{y}{r s_1^\#}
       (\log y)^{-2-\nu_1}.
\end{equation}

\paragraph{The prime conditions in an intermediate layer.}
Fix $2\le i\le b$ and $\sigma_i$.  The last slot, of index
$i-1$, is complete after this layer.  Let $Q_1,Q_2$ be its largest
prime factors on the two sides; both lie in
$[U_{i-1},Y_{i-1}]$.

Suppose first that $Q_1=Q_2=Q$.  Remove the unique occurrence of
$Q$ from $t_i^\#$, and call the result $t$.  For fixed allocations
of $t$, the completed primes are $AQ+1,A'Q+1$ with $A\ne A'$.
Lemma~\ref{lem:sieve} gives
\begin{equation}\label{eq:shared-layer-prime}
 \sum_Q\frac1Q\ll\frac{B_y^3}{(\log U_{i-1})^2}.
\end{equation}

Now suppose $Q_1\ne Q_2$ and remove both, obtaining $t$.
In addition to its two allocations, fix the position of $Q_1$ on
the second side and of $Q_2$ on the first side.  There are at most
$i^2$ choices.  If neither prime lies in the opposite last slot,
the two completed prime conditions are
$AQ_1+1$ and $A'Q_2+1$; two applications of
\eqref{eq:sieve-reciprocal} with $k=2$ give
\begin{equation}\label{eq:distinct-layer-primes}
 \sum_{Q_1,Q_2}\frac1{Q_1Q_2}
 \ll\frac{B_y^4}{(\log U_{i-1})^2}.
\end{equation}
If $Q_2$ lies in the last slot on the first side, then $Q_2<Q_1$.
The prime $Q_1$ cannot lie in the last slot on the second side,
whose largest prime is $Q_2$.  The completed prime conditions are
now $AQ_2Q_1+1$ and $A'Q_2+1$.  First sum over $Q_1$ conditional
on $Q_2$, using the forms $Q_1,AQ_2Q_1+1$.  Feasibility ensures
$AQ_2\le y$, so the inner reciprocal sum is uniformly
$O(B_y^2/\log U_{i-1})$.  Only after this estimate do we enlarge
the range of $Q_2$ and apply the two-form bound to
$Q_2,A'Q_2+1$.  This proves \eqref{eq:distinct-layer-primes} again.
The reversed incidence is treated by reversing the order of
summation.  Both incidences cannot occur simultaneously, as they
would imply both $Q_1<Q_2$ and $Q_2<Q_1$.

\paragraph{Counting the allocations.}
Put $\Delta_i=\nu_{i-1}-\nu_i$.  Normality bounds the number of
prime factors of an individual shift in $(Y_i,Y_{i-1}]$ by
$(\Delta_i+\delta)B_y$.  Indeed, if the shift ends before
$Y_{i-1}$, apply normality only up to its endpoint; if it ends
before $Y_i$, the count is zero.  In either case the error is at
most $\sqrt{\log_2 S\,B_y}=\delta B_y$.
Only $i$ shifts contribute to this layer, and removing the
distinguished primes decreases the count.  Thus
\[
 \Omega(t)\le I_i:=i(\Delta_i+\delta)B_y.
\]
Since $t$ is squarefree, each of its prime factors chooses one of
$i$ slots on each side, giving at most $i^{2\Omega(t)}$ dual
allocations.  Write
\[
 M_i=\sum_{Y_i<p\le Y_{i-1}}\frac1p.
\]
Mertens' estimate gives $M_i\le\Delta_i B_y+O(1)$, uniformly in
the cutoffs.  Since $\delta B_y\ge\sqrt{B_y}$, increasing the
absolute threshold $y_0$ ensures
$M_i\le(\Delta_i+\delta)B_y$.  Hence
\begin{align}
 \sum_t\frac{i^{2\Omega(t)}}t
 &\le\sum_{n\le I_i}\frac{(i^2M_i)^n}{n!}
 \le i^{I_i}\exp(iM_i)\notag\\
 &\le\exp\bigl((i\log i+i)(\Delta_i+\delta)B_y\bigr).
 \label{eq:allocation-entropy}
\end{align}
The middle inequality follows by writing $i^{2n}=i^ni^n$,
using $i^n\le i^{I_i}$, and extending the remaining exponential
series.  It is valid without any lower bound on
$\Delta_i/\delta$.

Combining \eqref{eq:shared-layer-prime},
\eqref{eq:distinct-layer-primes}, and
\eqref{eq:allocation-entropy}, and using $i\le b\le B_y$ to
absorb the $i^2$ position choices, gives
\begin{equation}\label{eq:intermediate-layer}
 \sum_{\substack{\sigma_{i-1}\text{ restricting}\\
                  \text{to }\sigma_i}}\frac1{t_i^\#}
 \le CB_y^6
 (\log y)^{-2\mu_{i-1}+(i\log i+i)(\nu_{i-1}-\nu_i+\delta)}.
\end{equation}
The squarefree hypothesis was used both to remove a single
occurrence of each distinguished prime and to obtain the
factorial denominator in \eqref{eq:allocation-entropy}.

\paragraph{Multiplying the layer estimates.}
Starting from \eqref{eq:top-layer}, sum over $\mathcal C_1$.
For each subsequent layer use the exact weighted identity
\[
 \sum_{\sigma_{i-1}\in\mathcal C_{i-1}}\frac1{s_{i-1}^\#}
 =\sum_{\sigma_i\in\mathcal C_i}\frac1{s_i^\#}
   \sum_{\sigma_{i-1}\text{ restricting to }\sigma_i}
          \frac1{t_i^\#}.
\]
All bounds in \eqref{eq:intermediate-layer} are uniform in the
lower configuration.  Their constants therefore multiply to at
most $C^bB_y^{6b}$.  If $c_i=i\log i+i$, so that $c_1=1$, the
power of $\log y$ accumulated in this procedure is
\begin{align*}
 &-2-\nu_1+
   \sum_{i=2}^b\{-2\mu_{i-1}+c_i(\nu_{i-1}-\nu_i+\delta)\}\\
 &\qquad=-2+\sum_{j=1}^{b-1}a_j\nu_j
     +2\sum_{j=1}^{b-1}(\nu_j-\mu_j)
     +\delta\sum_{i=2}^b c_i-c_b\nu_b.
\end{align*}
Here $c_{j+1}-c_j-2=a_j$, including $j=1$.  The identity also
holds for $b=1$.  Discard the nonpositive final term
$-c_b\nu_b$.  It remains to bound
$\sum_{\sigma_b\in\mathcal C_b}1/s_b^\#$.

\paragraph{The smooth parts and repeated factors.}
Write $s=s_b^\#/D$ and $L_b=\log_2Y_b$.  For a fixed $s$, its
allocation among the $b$ left slots has at most $b^{\Omega(s)}$
possibilities.  The allocation of $sD$ among the $b$ right slots
and the additional slot $D'$ has at most
$(b+1)^{\Omega(sD)}$ possibilities.  These bounds allow prime
powers: assigning labeled copies of a repeated prime to slots
overcounts every possible distribution of its exponent.

Each left shift has largest prime factor above $Y_b$.  Its number
of prime factors at most $Y_b$ is bounded by normality as follows:
\[
 \Omega(p_j-1,1,Y_b)
 \le 2\log_2 S+(L_b-\log_2 S)
       +\sqrt{L_b\log_2 S}\le3L_b.
\]
If $S=Y_b$, the first normality inequality alone gives the same
bound.  Thus $\Omega(s)\le3bL_b$.  The ordinary smooth-number
Euler product, including all powers, now gives
\begin{align*}
 \sum_{\sigma_b\in\mathcal C_b}\frac1{s_b^\#}
 &\le\frac{(b+1)^{\Omega(D)}}D
   [b(b+1)]^{3bL_b}
   \sum_{P^+(s)\le Y_b}\frac1s\\
 &=\frac{(b+1)^{\Omega(D)}}D
   [b(b+1)]^{3bL_b}
   \prod_{p\le Y_b}\left(1-\frac1p\right)^{-1}\\
 &\ll\frac{(b+1)^{\Omega(D)}}D
   (\log Y_b)^{1+3b\log(b(b+1))}.
\end{align*}
No squarefreeness is assumed below $Y_b$.  Combining this estimate
with the layer bounds proves \eqref{eq:comparison}, after increasing
the absolute constant $C$.
\end{proof}

\begin{remark}\label{rem:comparison-uniformity}
The explicit error in \eqref{eq:comparison-error} satisfies
\[
 E_b=2\sum_{j=1}^{b-1}(\nu_j-\mu_j)
       +O\bigl(\delta b^2\log(2b)\bigr),
\]
and the smooth-factor exponent is $O(b\log(2b))$.  Thus all
dependence on a growing $b$ is displayed in
\eqref{eq:comparison}.  No restriction such as
$r\le y^{1/10}$ or $D\le y^{1/100}$ is needed: the large top
part supplies the lower bound for the actual sieve endpoint.
\end{remark}
